% !TEX program = xelatex
% ============================================================
%  Arabic Business Plan Template — Nibras Center
%  خطة العمل للشركة الناشئة
%  Compile with: xelatex main.tex (twice)
% ============================================================

\documentclass[11pt,a4paper]{article}

\usepackage[a4paper,margin=2.5cm,top=2.5cm,bottom=2.5cm,headheight=15pt]{geometry}
\usepackage{array}
\usepackage{booktabs}
\usepackage{tabularx}
\usepackage{enumitem}
\usepackage{float}
\usepackage{titlesec}
\usepackage{fancyhdr}
\usepackage{setspace}
\usepackage[table]{xcolor}

\usepackage{bidi}
\usepackage[no-math]{fontspec}
\usepackage[quiet]{polyglossia}

\setmainfont[Scale=1]{NotoKufiArabic-Regular.ttf}
\newfontfamily\arabicfont[Script=Arabic,Scale=0.95]{NotoKufiArabic-Regular.ttf}
\newfontfamily\englishfont[Scale=0.95]{TeX Gyre Termes}

\setdefaultlanguage[calendar=gregorian,hijricorrection=1]{arabic}
\setotherlanguage[variant=british]{english}

\usepackage[breaklinks,hidelinks]{hyperref}
\hypersetup{pdftitle={خطة العمل},pdfauthor={الشركة}}

\titleformat{\section}{\Large\bfseries}{\thesection}{0.7em}{}
\titleformat{\subsection}{\large\bfseries}{\thesubsection}{0.6em}{}
\titlespacing*{\section}{0pt}{1.2ex plus 0.3ex}{0.8ex plus 0.2ex}

\pagestyle{fancy}
\fancyhf{}
\fancyhead[R]{\small\itshape خطة العمل}
\fancyhead[L]{\small اسم الشركة}
\fancyfoot[C]{\small\thepage}
\renewcommand{\headrulewidth}{0.3pt}

\newcolumntype{L}[1]{>{\raggedright\arraybackslash}p{#1}}
\newcolumntype{C}[1]{>{\centering\arraybackslash}p{#1}}

\setstretch{1.4}

\begin{document}

% ============================================================
% TITLE PAGE
% ============================================================
\thispagestyle{empty}
\begin{center}
\vspace*{3cm}
{\Huge\bfseries خطة عمل}\\[1em]
{\LARGE اسم الشركة}\\[2em]
{\Large شعار أو وصف مختصر}\\[3em]

\begin{tabular}{rl}
\textbf{المؤسسون:} & --- \\
\textbf{التاريخ:} & \today \\
\textbf{الإصدار:} & 1.0 \\
\textbf{السرية:} & خاص وسري \\
\end{tabular}
\end{center}

\newpage

% ============================================================
% 1. EXECUTIVE SUMMARY
% ============================================================
\section{الملخص التنفيذي}
\label{sec:summary}

% TODO: اكتب ملخصاً تنفيذياً (صفحة واحدة).
هذا نص تجريبي للملخص التنفيذي. يقدم نظرة عامة على الشركة والمنتج والسوق والفرصة والفريق والاحتياجات المالية.

\section{وصف الشركة}
\label{sec:company}

\subsection{النظرة العامة}
% TODO: صف الشركة.
\subsection{المهمة والرؤية}
\begin{itemize}
    \item \textbf{المهمة:} ---
    \item \textbf{الرؤية:} ---
\end{itemize}

\subsection{القيم الأساسية}
\begin{enumerate}
    \item ---
    \item ---
    \item ---
\end{enumerate}

% ============================================================
% 2. PROBLEM & SOLUTION
% ============================================================
\section{المشكلة والحل}
\label{sec:problem}

\subsection{المشكلة}
% TODO: صف المشكلة التي تحلها الشركة.

\subsection{الحل}
% TODO: صف المنتج أو الخدمة.

\subsection{الميزة التنافسية}
\begin{itemize}
    \item ---
    \item ---
\end{itemize}

% ============================================================
% 3. MARKET ANALYSIS
% ============================================================
\section{تحليل السوق}
\label{sec:market}

\subsection{حجم السوق}
\begin{table}[H]
\centering
\begin{tabular}{|l|c|c|}
\hline
\textbf{السوق} & \textbf{الحجم (\LR{USD})} & \textbf{معدل النمو} \\
\hline
السوق الكلي (\LR{TAM}) & --- & ---\% \\
السوق المتاح (\LR{SAM}) & --- & ---\% \\
السوق المستهدف (\LR{SOM}) & --- & ---\% \\
\hline
\end{tabular}
\end{table}

\subsection{العملاء المستهدفون}
% TODO: صف شريحة العملاء المستهدفة.

\subsection{المنافسون}
\begin{table}[H]
\centering
\renewcommand{\arraystretch}{1.3}
\begin{tabularx}{\textwidth}{L{3cm} X X X}
\toprule
\textbf{المنافس} & \textbf{نقاط القوة} & \textbf{نقاط الضعف} & \textbf{الاختلاف عنّا} \\
\midrule
--- & --- & --- & --- \\
--- & --- & --- & --- \\
--- & --- & --- & --- \\
\bottomrule
\end{tabularx}
\end{table}

% ============================================================
% 4. BUSINESS MODEL
% ============================================================
\section{نموذج العمل}
\label{sec:model}

\subsection{مصادر الإيراد}
\begin{enumerate}
    \item \textbf{المصدر 1:} --- (نسبة الإيراد: ---\%)
    \item \textbf{المصدر 2:} --- (نسبة الإيراد: ---\%)
\end{enumerate}

\subsection{التسعير}
\begin{table}[H]
\centering
\begin{tabular}{|l|c|c|}
\hline
\textbf{الباقة} & \textbf{السعر شهرياً} & \textbf{الميزات} \\
\hline
مجاني & 0 & --- \\
احترافي & --- \LR{USD} & --- \\
مؤسسات & --- \LR{USD} & --- \\
\hline
\end{tabular}
\end{table}

\subsection{التكاليف}
\begin{itemize}
    \item \textbf{تكاليف ثابتة:} ---
    \item \textbf{تكاليف متغيرة:} ---
    \item \textbf{تكلفة اكتساب العميل (\LR{CAC}):} ---
    \item \textbf{القيمة الدائمة للعميل (\LR{LTV}):} ---
    \item \textbf{نسبة \LR{LTV/CAC:} ---}
\end{itemize}

% ============================================================
% 5. GO-TO-MARKET
% ============================================================
\section{استراتيجية الوصول إلى السوق}
\label{sec:gtm}

\subsection{قنوات التسويق}
\begin{enumerate}
    \item \textbf{التسويق الرقمي:} ---
    \item \textbf{التسويق بالمحتوى:} ---
    \item \textbf{شراكات:} ---
    \item \textbf{مبيعات مباشرة:} ---
\end{enumerate}

\subsection{خطة الإطلاق}
\begin{itemize}
    \item \textbf{المرحلة 1 (شهر 1--3):} ---
    \item \textbf{المرحلة 2 (شهر 4--6):} ---
    \item \textbf{المرحلة 3 (شهر 7--12):} ---
\end{itemize}

% ============================================================
% 6. TEAM
% ============================================================
\section{الفريق}
\label{sec:team}

\begin{table}[H]
\centering
\renewcommand{\arraystretch}{1.4}
\begin{tabularx}{\textwidth}{L{4cm} L{4cm} X}
\toprule
\textbf{الاسم} & \textbf{المنصب} & \textbf{الخبرة} \\
\midrule
--- & \LR{CEO} & --- \\
--- & \LR{CTO} & --- \\
--- & \LR{CFO} & --- \\
\bottomrule
\end{tabularx}
\end{table}

% ============================================================
% 7. FINANCIAL PROJECTIONS
% ============================================================
\section{التوقعات المالية}
\label{sec:financials}

\begin{table}[H]
\centering
\renewcommand{\arraystretch}{1.4}
\begin{tabularx}{\textwidth}{L{4cm} C{2cm} C{2cm} C{2cm} C{2cm} C{2cm}}
\toprule
\textbf{البند} & \textbf{سنة 1} & \textbf{سنة 2} & \textbf{سنة 3} & \textbf{سنة 4} & \textbf{سنة 5} \\
\midrule
الإيرادات & 0 & 0 & 0 & 0 & 0 \\
تكلفة المبيعات & 0 & 0 & 0 & 0 & 0 \\
إجمالي الربح & 0 & 0 & 0 & 0 & 0 \\
المصاريف التشغيلية & 0 & 0 & 0 & 0 & 0 \\
صافي الربح & 0 & 0 & 0 & 0 & 0 \\
\bottomrule
\end{tabularx}
\end{table}

% ============================================================
% 8. FUNDING
% ============================================================
\section{الاحتياجات التمويلية}
\label{sec:funding}

\begin{itemize}
    \item \textbf{المبلغ المطلوب:} --- \LR{USD}
    \item \textbf{الجولة:} \LR{Seed / Series A}
    \item \textbf{التقييم قبل الاستثمار:} --- \LR{USD}
    \item \textbf{الأسهم المعروضة:} ---\%
\end{itemize}

\subsection{استخدام الأموال}
\begin{table}[H]
\centering
\begin{tabular}{|l|c|}
\hline
\textbf{البند} & \textbf{النسبة} \\
\hline
تطوير المنتج & 40\% \\
التسويق والمبيعات & 30\% \\
العمليات & 15\% \\
توظيف & 10\% \\
احتياطي & 5\% \\
\hline
\end{tabular}
\end{table}

% ============================================================
% 9. MILESTONES
% ============================================================
\section{المعالم الرئيسية}
\label{sec:milestones}

\begin{table}[H]
\centering
\renewcommand{\arraystretch}{1.4}
\begin{tabularx}{\textwidth}{C{1cm} L{6cm} C{3cm} X}
\toprule
\textbf{رقم} & \textbf{المعلم} & \textbf{التاريخ} & \textbf{معيار النجاح} \\
\midrule
1 & إطلاق النسخة التجريبية & --- & --- مستخدم \\
2 & الإطلاق الرسمي & --- & --- مستخدم \\
3 & الوصول إلى التعادل & --- & --- \LR{USD} إيراد \\
4 & جولة تمويل \LR{A} & --- & --- \LR{USD} \\
\bottomrule
\end{tabularx}
\end{table}

% ============================================================
% 10. RISKS
% ============================================================
\section{المخاطر}
\label{sec:risks}

\begin{table}[H]
\centering
\renewcommand{\arraystretch}{1.4}
\begin{tabularx}{\textwidth}{C{1cm} L{4cm} X}
\toprule
\textbf{رقم} & \textbf{المخاطرة} & \textbf{خطة المعالجة} \\
\midrule
1 & --- & --- \\
2 & --- & --- \\
3 & --- & --- \\
\bottomrule
\end{tabularx}
\end{table}

\end{document}
